% GENERATED FILE. Edit canonical lessons, then run npm run generate:books.
% canonical-source-hash: fnv1a64:ca824745b22ef07f
% canonical-lessons: GU-C448-hasya, GU-C448-hasavvu, GU-C448-smit, GU-C448-majak, GU-C448-cidavvu

\chapter{Laughing Together}
\label{ch:gu-laughing-together}
\hlchaptermodality{448}

\begin{chapteropening}
\textbf{By the end of this chapter:} \emph{I can say what makes us laugh.}

Complete the last lesson of chapter 448: I can say what makes us laugh.
\end{chapteropening}

\section[\texorpdfstring{\gu{હાસ્ય} (hāsya)}{હાસ્ય (hāsya)}]{\gu{હાસ્ય} (hāsya) — laughter}

\label{lesson:GU-C448-hasya}



\begin{quote}
Before the new one: say the Gujarati for to experience, then the Gujarati for to delight in.
\end{quote}

\subsection*{You'll want to know: \gu{હાસ્ય}}
\textbf{\gu{હાસ્ય}} — \emph{hāsya} — \textquotedblleft{}laughter\textquotedblright{}.

\begin{quote}
\gu{મીરાના} \gu{હાસ્યથી} \gu{આખું} \gu{ઘર} \gu{ખુશ} \gu{થઈ} \gu{ગયું}.
\end{quote}

\begin{quote}
\emph{Mīrānā hāsyathī ākhũ ghar khuś thaī gayũ.}
\end{quote}

\begin{quote}
Mira's laughter made the whole house happy.
\end{quote}

A neuter noun.

\begin{cousinweb}
\textbf{\gu{હાસ્ય}} (\emph{hāsya}) is a tatsama, Sanskrit \emph{hāsya}, \textquotedblleft{}laughter, mirth\textquotedblright{}, as Monier-Williams records it; it shares its root with \textbf{\gu{હસવું}} (\emph{hasvũ}), \textquotedblleft{}to laugh\textquotedblright{}.
\end{cousinweb}

\begin{grammarlens}[title={What you've built}]
One more word for describing what you lived and felt.
\end{grammarlens}

\subsection*{Your turn}
\begin{itemize}
\raggedright
  \item \emph{Say it:} \emph{hāsya}
  \item \emph{Say it:} \emph{hāsya}, once more
  \item \emph{Say it:} \emph{Mīrānā hāsyathī ākhũ ghar khuś thaī gayũ.} — Mira's laughter made the whole house happy.
  \item \emph{Recall:} say the Gujarati for to experience, then the Gujarati for to delight in, then say \emph{hāsya} again
\end{itemize}

\begin{tcolorbox}[breakable,colback=teal!4,colframe=teal!35!black,title={Before you move on}]
What does \gu{હાસ્ય} mean? (\textquotedblleft{}Laughter\textquotedblright{}.) Say it once more.
\end{tcolorbox}
\section[\texorpdfstring{\gu{હસાવવું} (hasāvvũ)}{હસાવવું (hasāvvũ)}]{\gu{હસાવવું} (hasāvvũ) — to make someone laugh}

\label{lesson:GU-C448-hasavvu}



\begin{quote}
Before the new one: say the Gujarati for to delight in, then the Gujarati for laughter.
\end{quote}

\subsection*{You'll want to know: \gu{હસાવવું}}
\textbf{\gu{હસાવવું}} — \emph{hasāvvũ} — \textquotedblleft{}to make someone laugh\textquotedblright{}.

\begin{quote}
\gu{અરુણની} \gu{વાર્તાઓ} \gu{બધાને} \gu{હસાવે} \gu{છે}.
\end{quote}

\begin{quote}
\emph{Aruṇnī vārtāo badhāne hasāve che.}
\end{quote}

\begin{quote}
Arun's stories make everyone laugh.
\end{quote}

Past, with \gu{તેણે} (he or she …): \gu{તેણે} \gu{બાળકને} \gu{હસાવ્યું}, \emph{teṇe bāḷakne hasāvyũ}.

\begin{cousinweb}
\textbf{\gu{હસાવવું}} (\emph{hasāvvũ}) is the causative of \textbf{\gu{હસવું}} (\emph{hasvũ}), \textquotedblleft{}to laugh\textquotedblright{}, which is inherited from Sanskrit \emph{hasati}, \textquotedblleft{}laughs\textquotedblright{} (CDIAL 14021).
\end{cousinweb}

\begin{grammarlens}[title={What you've built}]
One more word for describing what you lived and felt.
\end{grammarlens}

\subsection*{Your turn}
\begin{itemize}
\raggedright
  \item \emph{Say it:} \emph{hasāvvũ}
  \item \emph{Say it:} \emph{hasāvvũ}, once more
  \item \emph{Say it:} \emph{Aruṇnī vārtāo badhāne hasāve che.} — Arun's stories make everyone laugh.
  \item \emph{Recall:} say the Gujarati for to delight in, then the Gujarati for laughter, then say \emph{hasāvvũ} again
\end{itemize}

\begin{tcolorbox}[breakable,colback=teal!4,colframe=teal!35!black,title={Before you move on}]
What does \gu{હસાવવું} mean? (\textquotedblleft{}To make someone laugh\textquotedblright{}.) Say it once more.
\end{tcolorbox}
\section[\texorpdfstring{\gu{સ્મિત} (smit)}{સ્મિત (smit)}]{\gu{સ્મિત} (smit) — a smile}

\label{lesson:GU-C448-smit}



\begin{quote}
Before the new one: say the Gujarati for laughter, then the Gujarati for to make someone laugh.
\end{quote}

\subsection*{You'll want to know: \gu{સ્મિત}}
\textbf{\gu{સ્મિત}} — \emph{smit} — \textquotedblleft{}a smile\textquotedblright{}.

\begin{quote}
\gu{દાદીના} \gu{મોઢા} \gu{પર} \gu{મીઠું} \gu{સ્મિત} \gu{હતું}.
\end{quote}

\begin{quote}
\emph{dādīnā moḍhā par mīṭhũ smit hatũ.}
\end{quote}

\begin{quote}
There was a sweet smile on my grandmother's face.
\end{quote}

A neuter noun.

\begin{cousinweb}
\textbf{\gu{સ્મિત}} (\emph{smit}) is a tatsama, Sanskrit \emph{smita}, \textquotedblleft{}smiling\textquotedblright{}, as Monier-Williams records it; Gujarati uses it as the noun \textquotedblleft{}a smile\textquotedblright{}.
\end{cousinweb}

\begin{grammarlens}[title={What you've built}]
One more word for describing what you lived and felt.
\end{grammarlens}

\subsection*{Your turn}
\begin{itemize}
\raggedright
  \item \emph{Say it:} \emph{smit}
  \item \emph{Say it:} \emph{smit}, once more
  \item \emph{Say it:} \emph{dādīnā moḍhā par mīṭhũ smit hatũ.} — There was a sweet smile on my grandmother's face.
  \item \emph{Recall:} say the Gujarati for laughter, then the Gujarati for to make someone laugh, then say \emph{smit} again
\end{itemize}

\begin{tcolorbox}[breakable,colback=teal!4,colframe=teal!35!black,title={Before you move on}]
What does \gu{સ્મિત} mean? (\textquotedblleft{}A smile\textquotedblright{}.) Say it once more.
\end{tcolorbox}
\section[\texorpdfstring{\gu{મજાક} (majāk)}{મજાક (majāk)}]{\gu{મજાક} (majāk) — a joke; good-natured teasing}

\label{lesson:GU-C448-majak}



\begin{quote}
Before the new one: say the Gujarati for to make someone laugh, then the Gujarati for a smile.
\end{quote}

\subsection*{You'll want to know: \gu{મજાક}}
\textbf{\gu{મજાક}} — \emph{majāk} — \textquotedblleft{}a joke; good-natured teasing\textquotedblright{}.

\begin{quote}
\gu{અરુણે} \gu{મજાક} \gu{કરી} \gu{અને} \gu{બધા} \gu{હસ્યા}.
\end{quote}

\begin{quote}
\emph{Aruṇe majāk karī ane badhā hasyā.}
\end{quote}

\begin{quote}
Arun made a joke and everyone laughed.
\end{quote}

A feminine noun.

\begin{cousinweb}
\textbf{\gu{મજાક}} (\emph{majāk}) is usually traced to Arabic \emph{madhāq}, \textquotedblleft{}taste, relish\textquotedblright{}, which reached Gujarati through Persian, where it is said \emph{mazāq}; the sense \textquotedblleft{}a joke\textquotedblright{} is said to have grown up on the way.
\end{cousinweb}

\begin{grammarlens}[title={What you've built}]
One more word for describing what you lived and felt.
\end{grammarlens}

\subsection*{Your turn}
\begin{itemize}
\raggedright
  \item \emph{Say it:} \emph{majāk}
  \item \emph{Say it:} \emph{majāk}, once more
  \item \emph{Say it:} \emph{Aruṇe majāk karī ane badhā hasyā.} — Arun made a joke and everyone laughed.
  \item \emph{Recall:} say the Gujarati for to make someone laugh, then the Gujarati for a smile, then say \emph{majāk} again
\end{itemize}

\begin{tcolorbox}[breakable,colback=teal!4,colframe=teal!35!black,title={Before you move on}]
What does \gu{મજાક} mean? (\textquotedblleft{}A joke; good-natured teasing\textquotedblright{}.) Say it once more.
\end{tcolorbox}
\section[\texorpdfstring{\gu{ચીડવવું} (cīḍavvũ)}{ચીડવવું (cīḍavvũ)}]{\gu{ચીડવવું} (cīḍavvũ) — to tease}

\label{lesson:GU-C448-cidavvu}



\begin{quote}
Before the new one: say the Gujarati for a smile, then the Gujarati for a joke.
\end{quote}

\subsection*{You'll want to know: \gu{ચીડવવું}}
\textbf{\gu{ચીડવવું}} — \emph{cīḍavvũ} — \textquotedblleft{}to tease\textquotedblright{}.

\begin{quote}
\gu{મોટો} \gu{ભાઈ} \gu{નાની} \gu{બહેનને} \gu{રોજ} \gu{ચીડવે} \gu{છે}.
\end{quote}

\begin{quote}
\emph{moṭo bhāī nānī bahenne roj cīḍve che.}
\end{quote}

\begin{quote}
The big brother teases his little sister every day.
\end{quote}

Past, with \gu{તેણે} (he or she …): \gu{તેણે} \gu{બાળકને} \gu{ચીડવ્યું}, \emph{teṇe bāḷakne cīḍavyũ}.

\begin{cousinweb}
\textbf{\gu{ચીડવવું}} (\emph{cīḍavvũ}) is the causative partner of \textbf{\gu{ચિડાવું}} (\emph{ciḍāvũ}), \textquotedblleft{}to get irritated\textquotedblright{}: to tease is to make someone cross. Both go back to a reconstructed \emph{*ciḍ-}, \textquotedblleft{}be angry\textquotedblright{} (CDIAL 4794).
\end{cousinweb}

\begin{grammarlens}[title={What you've built}]
One more word for describing what you lived and felt.
\end{grammarlens}

\subsection*{Your turn}
\begin{itemize}
\raggedright
  \item \emph{Say it:} \emph{cīḍavvũ}
  \item \emph{Say it:} \emph{cīḍavvũ}, once more
  \item \emph{Say it:} \emph{moṭo bhāī nānī bahenne roj cīḍve che.} — The big brother teases his little sister every day.
  \item \emph{Recall:} say the Gujarati for a smile, then the Gujarati for a joke, then say \emph{cīḍavvũ} again
\end{itemize}

\begin{tcolorbox}[breakable,colback=teal!4,colframe=teal!35!black,title={Before you move on}]
What does \gu{ચીડવવું} mean? (\textquotedblleft{}To tease\textquotedblright{}.) Say it once more.
\end{tcolorbox}
